% Abstract v4 - attractor-first thesis. Physics as discovery tool, not claim.
% Honest about margin-sensitivity; structural attractor as the headline.
\begin{abstract}
Load balancing on a graph is usually framed as an algorithm-design problem: a
better router yields a better load distribution. We report evidence that, over
a wide range of topologies, the load distribution is instead largely fixed by
the structure of the graph and the traffic demand, leaving any reasonable
congestion-aware router little freedom---and we identify a topological metric
that predicts how much freedom remains. We reached this through an unusual
route. Treating a packet as a particle under a classical force, we built three
independent routing cores---a Newtonian loss-reaction, an Archimedean buoyancy,
and a Pascalian pressure diffusion---sharing a search engine but no mechanism.
The three produce near-identical per-edge load distributions (cosine $0.99$),
significantly tighter than blind controls such as ECMP ($0.87$), and reach
statistical equivalence (TOST) with the engineered balancers CONGA and DRILL on
most of fifteen topologies. A sensitivity analysis is candid about the limits of
that equivalence: at a strict margin ($0.5$ percentage points of drop rate) only
a few topologies remain equivalent, and they are precisely the ones where the
structural attractor is strongest. We argue, via a cut-based structural
argument, that this is no coincidence: the demand fixes the flow across every
cut, the topology's cut structure pins most of the load distribution, and any
congestion-aware rule---physical or engineered---shapes the small residual
identically. A purely structural ratio, tube/sp (room to manoeuvre among
near-shortest paths, computable before deployment), predicts where the residual
is large enough for routers to differ; the correlation is robust to leave-one-out
and strengthens when the loosest-fitting topologies are removed ($r$ rising from
$0.78$ to $0.89$). The physics, in the end, was a discovery instrument rather
than the result: three classical analogies converged because the problem, not
the analogy, drove them to the same place. We report our negative results, an
audit retracting an earlier inflated advantage, and the margin sensitivity in
full.
\end{abstract}
